%%%%%%%%%%  scan idx 182  |  folio 169  %%%%%%%%%%
\item[(i)] \emph{Les
ouverts $S_f$, pour $f\in\Gamma(S,L_i^{\otimes n})$, $i\in I$ et $n\in\mathbb N$, forment une base de la topologie de} $S$.
\item[(i$'$)] \emph{Ceux des $S_f$ qui sont affines recouvrent} $S$.
\item[(ii)] \emph{Pour tout $\mathcal O_S$-Module $F$ quasi-cohérent et de type fini, il existe une famille $(r_i)_{i\in I}$ (resp. $(n_i)_{i\in I}$) d'entiers $\ge0$ (resp. $>0$), et un épimorphisme}
\[
\coprod_{i\in I}\mathcal O_S^{r_i}\otimes L_i^{\otimes-n_i}\longrightarrow F.
\]
\item[(iii)] \emph{Pour tout $F\in\Ob\Qcoh(S)$, on a}
\[
F=\sum_{(i,n)\in I\times\mathbb N^+}F_{(i,n)}\otimes L_i^{\otimes-n}
\]
(\emph{somme (non directe) de sous-faisceaux de} $F$).
\item[(iii$'$)] \emph{Même condition que (iii), mais $F$ étant un Idéal quasi-cohérent de} $\mathcal O_S$.
\end{enumerate}

\emph{Preuve.} Elle consiste en une simple paraphrase de (loc. cit.) ; nous la donnons, brièvement, pour la commodité du lecteur.

(i)$\Rightarrow$(i$'$) : Pour $x\in S$, on choisit un ouvert affine $U$ contenant $x$. Un ouvert $S_f$ tel que $x\in S_f\subset U$, avec $f\in\Gamma(S,L_i^{\otimes n})$ ($n\ge0$), est nécessairement affine, comme il résulte du

\textbf{Lemme 2.2.3.1.} (EGA II 5.5.8). \emph{Soient $X$ un schéma, $L$ un $\mathcal O_X$-Module inversible $f\in\Gamma(X,L)$. L'immersion canonique $i:X_f\to X$ est affine}.

\emph{Preuve.} C'est trivial : la question est locale sur $X$, donc on peut supposer $X$ affine et $L\simeq\mathcal O_X$, et l'on gagne.

(i$'$)$\Rightarrow$(ii) : Soit un recouvrement fini de $S$ par des ouverts affines $S_{f_{ij}}$, où $f_{ij}\in\Gamma(S,L_i^{\otimes m_{ij}})$ ; posons $S_i=\bigcup_j S_{f_{ij}}$. Quitte à remplacer les $f_{ij}$ par des puissances convenables, on peut supposer que les $m_{ij}$ correspondant à un même indice $i$ sont égaux à un entier $m_i$. Le faisceau 
